\documentclass[a4paper,10pt]{report}\input{../0.Préambule/Préambule_cours_prof}\begin{document}

\newpage
\section*{Enchainement d'opérations}
\addcontentsline{toc}{section}{Enchainement d'opérations}

\subsection*{Calculer une expression sans parenthèses}
\addcontentsline{toc}{subsection}{Calculer une expression sans parenthèses}

\begin{exer1}
Effectue les calculs suivants en soulignant le calcul en cours.

\begin{enumerate}[font=\color{exer1}]
\begin{tabularx}{\linewidth}{*{3}{>{\arraybackslash}X}}
\item $A=14-5+3$ &
\item $B=14+5-3$ &
\item $J=15+5+3$ \\
\item $E=24+19-5$ &
\item $C=24-19-5$ &
\item $T=24-5-19$ \\
\item $I=3\times 2\times 11$ &
\item $O=2\times 4 \div 4$ &
\item $N=15 \div 5 \times 8$ \\
\end{tabularx}
\end{enumerate}
\end{exer1}

\begin{exer1}
Entoure le signe opératoire de l'opération prioritaire. (Il peut y en avoir plusieurs)

\begin{enumerate}[font=\color{exer1}]
\begin{tabularx}{\linewidth}{*{4}{>{\arraybackslash}X}}
\item $O = 252 + 21 \times 41$ &
\item $X = 6,3 - 2,1 \div 7$ &
\item $Y = 3 + 0,3  \times -3$ &
\item $D = 2 \times 2 - 2 \div 2$ \\
\item $A = 17 - 15 \div 3 + 1$ &
\item $B = 50 + 3 + 2 \times 10$ &
\item $L = 0,204 \times 99 - 5,4$ &
\item $E =9 + 12 \times 11 \div 8$ \\
\end{tabularx}
\end{enumerate}
\end{exer1}

\begin{exer1}
Effectue les calculs suivants en soulignant le(s) calcul(s) en cours.

\begin{enumerate}[font=\color{exer1}]
\begin{tabularx}{\linewidth}{*{3}{>{\arraybackslash}X}}
\item $J = 24 + 3 \times 7$ &
\item $O = 720 \div 9 + 4$ &
\item $C = 15 \div 5 - 2$ \\
\item $K = 20 - 0,1 \times 38$ &
\item $E = 60 - 14 + 5 \times 3 + 2$ &
\item $Y = 8 \times 3 - 5 \times 4 \times 0,2$ \\
\end{tabularx}
\end{enumerate}
\end{exer1}

\begin{exer1}
Avec la calculatrice, calcule les expressions suivantes sans noter les résultats intermédiaires.

\begin{enumerate}[font=\color{exer1}]
\begin{tabularx}{\linewidth}{*{3}{>{\arraybackslash}X}}
\item $C = 43,21 - 17,03 + 132,11 - 61,45$ &&
\item $H = 3,15 \times 5,2 \times 2,5$ \\
\item $Y = 6,21 \times 3 + 4,01 \times 1,5$ &
\item $M = 3,15 \div 0,5 \div 2,5$ &
\item $E = 9,21 \div 3 - 4,02 \div 1,5$ \\
\end{tabularx}
\end{enumerate}
\end{exer1}

\begin{exer1}
Complète avec les signes $+$, $-$, $\times$ et $\div$ pour que les égalités soient vraies.

\begin{enumerate}[font=\color{exer1}]
\begin{tabularx}{\linewidth}{*{4}{>{\arraybackslash}X}}
\item $5 ~\dots~ 8 ~\dots~ 2 = 20$ &
\item $7 ~\dots~ 5 ~\dots~ 5 = 6$ &
\item $8 ~\dots~ 6 ~\dots~ 2 = 24$ &
\item $8 ~\dots~ 2 ~\dots~ 81 = 324$ \\
\end{tabularx}
\end{enumerate}
\end{exer1}

\begin{exer1} Avec un ordre de grandeur
\begin{enumerate}[font=\color{exer1}]
\item Détermine un ordre de grandeur de chacun des nombres suivants.
\begin{enumerate}[font=\color{exer1}]
\begin{tabularx}{\linewidth}{*{3}{>{\arraybackslash}X}}
\item $C = 3,15 \times 95,2 - 4,22$ &
\item $H = \np{40129,5} + 103,2 \times 98,017$ &
\item $O = 103,7272 \div 9,86 \times 489,7$ \\
\item $Y = \np{8109,8} - 3,204 \times 324,48$ &
\item $E = 17,025 + 49,892 \times \np{2015,8}$ &
\item $R = \np{9026,9} \div 101,19 - 0,78$ \\
\end{tabularx}
\end{enumerate}
\item Avec ta calculatrice, trouve la valeur exacte de chacun de ces nombres afin de vérifier.
\end{enumerate}
\end{exer1}

\subsection*{Calculer une expression avec parenthèses}
\addcontentsline{toc}{subsection}{Calculer une expression avec parenthèses}

\begin{exer1}
Entoure le signe opératoire de l'opération prioritaire. (Il peut y en avoir plusieurs)

\begin{enumerate}[font=\color{exer1}]
\begin{tabularx}{\linewidth}{*{4}{>{\arraybackslash}X}}
\item $K = (6,2 - 0,1) \div 10$ &
\item $Y = 238 - 4 \times (13 + 27)$ &
\item $S = 5 + (2,8 + 6 \times 1,2)$ &
\item $T = 90 - (2 \times 7 - 7) \times 6$ \\
\item $I = 9 \div 3 + (15 - 6 \div 3)$ &
\item $Q = 34 - ( 104 \div 52 \times 6)$ &
\item $U = (84 - 1) \div (5 + 0,4)$ &
\item $E = 3 \times \left( (1+2) \times 4 - 2 \right)$ \\
\end{tabularx}
\end{enumerate}
\end{exer1}

\begin{exer1}
Effectue les calculs en soulignant le calcul en cours.

\begin{enumerate}[font=\color{exer1}]
\begin{tabularx}{\linewidth}{*{3}{>{\arraybackslash}X}}
\item $A = 25 - (8 - 3) + 1$ &
\item $Z = 25 - 8 - (3 + 1)$ &
\item $Y = 25 - (8 - 3 + 1)$ \\
\item $M = 18 - \left( 4 \times (5 - 3) + 2 \right)$ &
\item $E = 24 \div \left(8 - (3 + 1) \right)$ &
\item $S = \left( 2 + 0,1 \times (5 + 3) \right) \div 4$ \\
\end{tabularx}
\end{enumerate}
\end{exer1}

\begin{exer1}
Observe puis calcule astucieusement les expressions suivantes.

\begin{enumerate}[font=\color{exer1}]
\begin{tabularx}{\linewidth}{*{2}{>{\arraybackslash}X}}
\item $W = (52 \times 321 - 18 \times 25) \times (2 \times 31 -62)$ &\\
\item $O = (78 + 7 \times 27) \div (78 + 7 \times 27))$ &
\item $K = 0,4 \times 0,27 \times 250$ \\
\end{tabularx}
\end{enumerate}
\end{exer1}

\begin{exer1}
Avec la calculatrice, calcule les expressions suivantes sans noter les résultats intermédiaires.

\begin{enumerate}[font=\color{exer1}]
\begin{tabularx}{\linewidth}{*{2}{>{\arraybackslash}X}}
\item $E = 54,2 - (8,72 - 5,21)$ &
\item $R = 7,2 \times (15,7 + 0,51) \times 3,5$ \\
\item $G = \left( ( 19,01 - 7,5) \times 2 - 13,02 \right) \times 2,3$ &
\item $S = \left( ( 20,52 + 7,5) \times 2 \right) \times (13 - 2,3)$ \\
\end{tabularx}
\end{enumerate}
\end{exer1}

\begin{exer1}
Réécris chaque expressions en supprimant, quand c'est possible, les parenthèses inutiles.

\begin{enumerate}[font=\color{exer1}]
\begin{tabularx}{\linewidth}{*{3}{>{\arraybackslash}X}}
\item $A = 21 - ( 8 \times 4)$ &
\item $B = 21 \times (8 - 4)$ &
\item $R = 21 - (8 - 4)$ \\
\item $I = (21 \times 8) - 4$ &
\item $E = (21 + 8 - 1) \div 4$ &
\item $Z = 21 - \left(8 - (4 \times 2) \right)$ \\
\end{tabularx}
\end{enumerate}
\end{exer1}

\begin{exer1}
Place les parenthèses pour que les égalités soient vraies et vérifie tes réponses.

\begin{enumerate}[font=\color{exer1}]
\begin{tabularx}{\linewidth}{*{4}{>{\arraybackslash}X}}
\item $4 \times 2 + 9 = 44$ &
\item $15 - 3 \times 2 = 24$ &
\item $5 + 5 \times 5 - 5 = 0$ &
\item $2 \times 5 - 2 \times 4 + 1 = 30$ \\
\end{tabularx}
\end{enumerate}
\end{exer1}

\begin{exer1}
Calcule les expressions suivantes.

\vspace{-0.75cm}
\begin{enumerate}[font=\color{exer1}]
\begin{tabularx}{\linewidth}{*{3}{>{\arraybackslash}X}}
&&\item $K = 35 - \left( 4 \times (5 + 2) - 7 \right)$ \\
\item $A = 12 \times \left( 32 - (4 + 7) \times 2 \right)$ &
\item $M = (1 + 7) \times \left( 11 - (2 + 3) \right)$ &
\item $I = 12 + \left( (120 - 20) - 2 \times 4 \times 5 \right)$ \\
\end{tabularx}
\end{enumerate}
\end{exer1}

\subsection*{Expressions avec un quotient}
\addcontentsline{toc}{subsection}{Expressions avec un quotient}

\begin{exer1}
Calcul en détaillant les étapes.

\vspace{-1cm}
\begin{enumerate}[font=\color{exer1}]
\begin{tabularx}{\linewidth}{*{4}{>{\arraybackslash}X}}
&&\item $R = \dfrac{5+3}{2}$ &
\item $A = \dfrac{9}{4-1}$ \\
\end{tabularx}
\end{enumerate}
\end{exer1}

\begin{exer1}
Ecris les expressions sous la forme d'un calcul en ligne (N'oublie pas les parenthèses !)

\begin{enumerate}[font=\color{exer1}]
\begin{tabular}{m{1.5cm}m{2cm}m{2.5cm}m{1.5cm}m{3cm}m{3cm}}
\item $8 + \dfrac{5}{4}$ &
\item $\dfrac{17 - 15}{3 + 2}$ &
\item $17 - \dfrac{15}{3} + 2$ &
\item $\dfrac{8}{5 + 4}$ &
\item $17 \times \dfrac{15 \times 4}{3 - 2} + 2 \times 8$ &
\item $\dfrac{15 + 4}{13 - 3} - 0,3 \times 10$ \\
\end{tabular}
\end{enumerate}
\end{exer1}

\begin{exer1}
Calcule chacune des expressions suivantes.

\begin{enumerate}[font=\color{exer1}]
\begin{tabularx}{\linewidth}{*{3}{>{\arraybackslash}X}}
\item $A = \dfrac{81}{9} \times 5 - 1$ &
\item $V = \dfrac{45,5}{2} \times 3 - 1$ &
\item $I = \dfrac{27}{2 \times 3} - 1$ \\
\item $S = \dfrac{17 - 5}{3} + 2$ &
\item $E = 7 \times \dfrac{15 \times 4}{3 - 2} + 2 \times 8$ &
\item $Z = \dfrac{13 \times (4 + 7) - 5}{13 - (2 \times 4 + 3}$ \\
\end{tabularx}
\end{enumerate}
\end{exer1}

\subsection*{Résoudre un problème}
\addcontentsline{toc}{subsection}{Résoudre un problème}

\begin{exer2}
Voici un programme de calcul 

Applique ce programme à chacun des nombres en écrivant \\une seule expression permettant de trouver le résultat puis calcule :

\begin{enumerate}[font=\color{exer2}]
\begin{tabularx}{\linewidth}{*{4}{>{\arraybackslash}X}}
\item $2,25$ &
\item $4,7$ &&\\
\end{tabularx}
\end{enumerate}

\vspace{-38mm}
\begin{flushright}
\arrayrulecolor{black}
\renewcommand{\arraystretch}{0.1}
\begin{tabular}{|p{5cm}|}
\hline
\begin{verbatim} Choisis un nombre \end{verbatim}
\begin{verbatim} Ajoute 3 \end{verbatim}
\begin{verbatim} Multiplie par 5 \end{verbatim}
\begin{verbatim} Retire 9 \end{verbatim}
\\\hline
\end{tabular}
\end{flushright}
\end{exer2}

\begin{exer2}
Voici quatre nombres : \quad\quad $12,5$ \quad\quad $8$ \quad\quad $6,5$ \quad\quad $1,2$

\noindent Pour chaque question, tu ne peux utiliser qu'une fois exactement les quatre nombres, l'addition, la soustraction et la multiplication. Toutefois, tu peux placer des parenthèses. Le résultat doit être positif.

\noindent Écris l'expression qui donne :
\begin{enumerate}[font=\color{exer2}]
\begin{tabularx}{\linewidth}{*{2}{>{\arraybackslash}X}}
\item le plus grand résultat possible &
\item le plus petit résultat possible\\
\end{tabularx}
\end{enumerate}
\end{exer2}

\begin{exer2}
Maya achète $5$ pots de confitures à $1,80$\euro pièce et $12$ baguettes de pain à $0,70$\euro pièce. 

\noindent Quel est le prix total qu'elle doit payer ?
\end{exer2}

\begin{exer3}
Une famille de $4$ personnes produit en moyenne $1,42$t de déchets par an. La population française est estimé à environ $67,1$ millions de personnes. 

\noindent Calcul, en kg, la masse de déchets produite de France chaque année. On rappelle que $1$t = $\np{1000}$kg.
\end{exer3}
\end{document}